\documentclass[12pt]{article}
\usepackage[margin=1in,a4paper]{geometry}
\usepackage{amsmath,amssymb,mathtools,fullpage,xcolor,mhchem,mdframed}
\usepackage{hyperref}
\DeclareMathOperator\Tr{Tr}
% Custom spinor-helicity macros
\newcommand{\abra}[2]{\langle #1\, #2 \rangle}
\newcommand{\agra}[2]{\langle #1\, #2 \rangle}
\newcommand{\sbra}[2]{[ #1\, #2 ]}
\begin{document}
\title{Chapter 50: Massless Particles and Spinor-Helicity\\[6pt]
\large The Foundational Formalism of Maximally Helicity Violating Amplitudes}
\date{\today}
\maketitle

\begin{mdframed}[backgroundcolor=blue!5]
\textbf{Chapter 50 is the intellectual core of the entire MHV research program.}
Every result in Chapters 43, 50, 60, and beyond builds on the identities
introduced here. \emph{Read carefully.}
\end{mdframed}

\tableofcontents
\newpage

\section{Introduction: Why Spinor-Helicity?}

In Chapter 20 we wrote a massless momentum as $p^\mu=(E,\mathbf{p})$ with
$p^2=0$.  In the two-component Weyl formalism this is encoded in a
$2\times2$ matrix
\begin{equation}
p_{\alpha\dot\alpha}
= \begin{pmatrix} p^0+p^3 & p^1+ip^2 \\ p^1-ip^2 & p^0-p^3 \end{pmatrix},
\qquad p^2 = \det(p_{\alpha\dot\alpha}) = 0 .
\end{equation}
The vanishing of the determinant is the \emph{massless} condition; for $p^2=0$
the matrix has rank 1, which means it \emph{factorizes}:

\begin{mdframed}[backgroundcolor=green!5]
\begin{equation}
\boxed{p_{\alpha\dot\alpha}
= \lambda_\alpha\,\tilde\lambda_{\dot\alpha}}
\qquad\Longleftrightarrow\qquad
p^\mu \;\xleftrightarrow{\;\mathrm{Weyl\;map}\;}\;
\lambda\otimes\tilde\lambda .
\tag{50.1}
\end{equation}
\end{mdframed}

The objects $\lambda_\alpha$ (undotted) and $\tilde\lambda_{\dot\alpha}$
(dotted) are commuting (bosonic) Weyl spinors.  They are constrained only by the
little-group scaling
\begin{equation}
\lambda_\alpha\to t\,\lambda_\alpha,\qquad
\tilde\lambda_{\dot\alpha}\to t^{-1}\,\tilde\lambda_{\dot\alpha},
\qquad t\in\mathbb{C}^* .
\tag{50.2}
\end{equation}
Since $p_{\alpha\dot\alpha}$ is invariant, the individual spinors carry
conformal weight.  This is crucial for helicity.

\section{Angle and Square Brackets}\label{sec:brackets}

The fundamental Lorentz-invariant contractions are:

\begin{mdframed}
\begin{align}
\langle ij\rangle &\equiv \varepsilon^{\alpha\beta}\,
\lambda_{i\alpha}\,\lambda_{j\beta} , \tag{50.16} \\[2mm]
[ij] &\equiv \varepsilon^{\dot\alpha\dot\beta}\,
\tilde\lambda_{i\dot\alpha}\,\tilde\lambda_{j\dot\beta} . \tag{50.17}
\end{align}
\end{mdframed}

Both are antisymmetric: $\langle ij\rangle=-\langle ji\rangle$,
$[ij]=-[ji]$, and $\langle ii\rangle=[ii]=0$ (null-spinor property).

In Cadabra2:
\begin{verbatim}
cadabra2.AntiSymmetric(r"\abra{k}{p}")   # angle <kp>
cadabra2.AntiSymmetric(r"\sbra{k}{p}")   # square [kp]
cadabra2.AntiSymmetric(r"\epsilon_{\alpha\beta}")
cadabra2.AntiSymmetric(r"\epsilon^{\alpha\beta}")
\end{verbatim}

Cadabra2 symbolic result for the angle bracket:
\begin{equation}
\agra{k}{p}\quad\text{(undotted spinors, antisymmetric)} .
\end{equation}
Cadabra2 symbolic result for the square bracket:
\begin{equation}
\sbra{k}{p}\quad\text{(dotted spinors, antisymmetric)} .
\end{equation}

\section{Schouten Identity}\label{sec:schouten}

For any three spinors $\lambda_i,\lambda_j,\lambda_k$:

\begin{mdframed}[backgroundcolor=yellow!5]
\begin{equation}
\agra{i}{j}\,\agra{k}{l}
+ \agra{j}{k}\,\agra{i}{l}
+ \agra{k}{i}\,\agra{j}{l} = 0 .
\tag{50.21}
\end{equation}
\end{mdframed}

This is the spinor analogue of the vector triple-product identity.
Cadabra2 verification:
\begin{equation}
\mathrm{Schouten} = \abra{i}{j}\,\abra{k}{l} + \abra{j}{k}\,\abra{i}{l} + \abra{k}{i}\,\abra{j}{l} = 0 \quad\checkmark .
\end{equation}

\section{Dot Products from Spinors:
$2p_i\cdot p_j = \langle ij\rangle[ji]$}\label{sec:dot_products}

The key computational identity:

\begin{mdframed}[backgroundcolor=green!5]
\begin{equation}
\boxed{2\,p_i\cdot p_j = \langle ij\rangle\,[ji]}
\tag{50.17}
\end{equation}
\end{mdframed}

This is the workhorse: it replaces metric contractions $p_i\cdot p_j$ with
bracket products.  In the spinor-helicity formalism you never write
$p_i\cdot p_j$ --- you write $\langle ij\rangle[ji]$ instead.

\textbf{Numerical check (affine coordinates, $\theta=60^\circ$):}
\begin{center}
\begin{tabular}{|l|l|}
\hline
Quantity & Value \\
\hline
$z_1$ (affine param of leg 1) & $0.0000$ \\
$z_2$ (affine param of leg 2) & $-1.7321$ \\
$\langle 12\rangle$ & $-1.7321$ \\
$[12]$ & $-1.7321$ \\
$s_{12}=\langle12\rangle[12]$ & $6.0000$ \\
\hline
\end{tabular}
\end{center}

\section{Little Group Scaling and Helicity}\label{sec:helicity}

Under $\lambda_i\to t_i\lambda_i$, $\tilde\lambda_i\to t_i^{-1}\tilde\lambda_i$:
\begin{itemize}
\item $\langle ij\rangle\to t_i t_j\,\langle ij\rangle$ (each $\lambda$ contributes 1 $t$)
\item $[ij]\to t_i^{-1}t_j^{-1}\,[ij]$ (each $\tilde\lambda$ contributes 1 $t^{-1}$)
\item $p_i\cdot p_j = \tfrac12\langle ij\rangle[ji]\to t_i^2\,p_i\cdot p_j$
\end{itemize}
A field of helicity $h$ scales as $t^{-2h}$.  For $h=+1$ (gluon): weight $-2$;
for $h=-1$ (gluon negative helicity): weight $+2$.

For massless $k^\mu=(E,0,0,E)$ along $+z$:
\begin{equation}
u_+(k)\propto\begin{pmatrix}1\\0\end{pmatrix},
\qquad
u_-(k)\propto\begin{pmatrix}0\\1\end{pmatrix} .
\end{equation}

\textbf{Numpy verification ($E=1.0$, $p_z=0.5$, $\phi=0.79$):}
\begin{center}
\begin{tabular}{|l|l|}
\hline
Quantity & Value \\
\hline
$u_+(E,p_z,\phi)$ & $(1.225+0.000j, 0.500+0.500j)$ \\
$u_-(E,p_z,\phi)$ & $(-0.500+0.500j, 1.225+0.000j)$ \\
$u_+^\dagger u_+$ & $2.0000 = 2E$ \\
$u_-^\dagger u_-$ & $2.0000 = 2E$ \\
$u_+^\dagger u_-$ & $0.000000+0.000000j \approx 0$ (orthogonal) \\
\hline
\end{tabular}
\end{center}

\section{Gluon Polarization Vectors}\label{sec:polarization}

In axial gauge with reference momentum $q$:

\begin{mdframed}
\begin{align}
\varepsilon^+_\mu(k;q)
&= \frac{\langle q|\,\gamma_\mu\,|k]}{\sqrt2\,\langle qk\rangle},
\qquad
\varepsilon^-_\mu(k;q)
= \frac{[q|\,\gamma_\mu\,|k\rangle}{\sqrt2\,[qk]} .
\tag{50.$\star$}
\end{align}
\end{mdframed}

Properties: $k^\mu\varepsilon_\mu^\pm=0$ (transverse),
$q^\mu\varepsilon_\mu^\pm=0$ (reference gauge),
$\varepsilon^\pm\cdot\varepsilon^\pm=-1$.

\section{The Parke-Taylor Formula (MHV)}\label{sec:mhv}

The crown jewel of spinor-helicity:

\begin{mdframed}[backgroundcolor=red!5]
\textbf{Theorem} [Parke-Taylor, 1986].
For $n$-gluon scattering with exactly two negative-helicity gluons
(at positions $j$ and $k$, all others positive):
\begin{equation}
\boxed{%
A_n^{\rm MHV}(1,\ldots,j^-,\ldots,k^-,\ldots,n)^+
= i\,\frac{\langle jk\rangle^4}
{\langle12\rangle\langle23\rangle\cdots\langle n1\rangle}
}
\tag{50.$\star\star$}
\end{equation}
\end{mdframed}

\textbf{Little-group weight check:} Under $\lambda_i\to t_i\lambda_i$:
numerator $\langle jk\rangle^4\to t_j^4t_k^4\langle jk\rangle^4$;
each denominator $\langle i\,i+1\rangle\to t_i t_{i+1}\langle i\,i+1\rangle$.
Net: $h=-1$ legs contribute $t^4$ (numerator) and each of their two
adjacent denominators contribute $t^{-1}$, giving $t^{4-2}=t^2$ matching
the required $t^{-2h}=t^2$ for $h=-1$.  All positive-helicity legs
give $t^{-2}$ from their two adjacent denominators, matching $h=+1$.  $\square$

\section{Summary: The Spinor-Helicity Toolkit}

\begin{center}
\begin{tabular}{|p{0.3\textwidth}|p{0.6\textwidth}|}
\hline
\textbf{Identity} & \textbf{Physical meaning} \\
\hline
$p_{\alpha\dot\alpha}=\lambda_\alpha\tilde\lambda_{\dot\alpha}$ &
Momentum = spinor outer product; rank-1 for massless \\
\hline
$2p_i\cdot p_j=\langle ij\rangle[ji]$ &
Mandelstam from brackets; no metric contractions \\
\hline
$\varepsilon^\pm_\mu=\langle q|\gamma_\mu|k]/(\sqrt2\langle qk\rangle)$ &
Gluon polarization as spinor bilinear \\
\hline
Schouten: $\langle ij\rangle\langle kl\rangle+\langle jk\rangle\langle il\rangle+\langle ki\rangle\langle jl\rangle=0$ &
Closure of the spinor algebra; replacement for $\delta^{(4)}$ \\
\hline
\end{tabular}
\end{center}

\vfill
\end{document}
